\chapter{Formulation complète du programme}
\label{ann:formulation}
% BUDGET : 2 pages. Version 2026-09-02 : l'annexe ne contient plus que des tableaux, la prose
% qui portait de l'information ayant ete convertie en lignes de table.
% La correspondance equation par equation est en annexe B.

\section{Ensembles, paramètres, variables}

\begin{table}[H]
  \centering
  \caption{Ensembles, variable de décision et paramètres du programme, avec l'unité ou le code
  de chacun et l'origine de la donnée.}
  \label{tab:ann-ensembles}
  \scriptsize
  \setlength{\tabcolsep}{4pt}
  \begin{tabularx}{\textwidth}{@{}llX@{}}
    \toprule
    \textbf{Symbole} & \textbf{Unité ou code} & \textbf{Contenu et origine} \\
    \midrule
    $P$ & \texttt{PLOTS} & parcelles du \gls{rpg} 2017, identifiants synthétiques
      \texttt{P1..Pn} \\
    $C$ & \texttt{CROPS} & activités \enquote{culture $\times$ itinéraire technique},
      84 codes fins \\
    $F$ & \texttt{FARMS} & exploitations, $f(p)$ désignant celle de la parcelle $p$ \\
    $E$ & \texttt{PAIRS} & couples éligibles $(p,c)$, seuls déclarés au solveur \\
    $Y_{p,c}$ & \texttt{model.Y} & variable binaire d'affectation, $(p,c) \in E$ \\
    \midrule
    $s_p$ & ha & surface de la parcelle, table parcellaire \\
    $m_c$ & \eur/ha & marge brute unitaire calculée à partir des produits, des charges et des
      aides \gls{posei}. Dépend de l'année économique choisie en configuration \\
    $a_f$ & sans & aversion au risque, déduite du type d'exploitation, lui-même déduit du
      groupe de cultures dominant. Huit valeurs, Table~2 de \textcite{chopin2015} \\
    $v_c$ & sans & fraction de marge perdue en cas d'aléa. Ce n'est pas une variance. Aucune
      covariance entre cultures n'intervient et la pénalité est linéaire en surface \\
    $h_c$ & h/ha & besoin en travail de l'itinéraire technique, Table~1 de l'article \\
    $H_f$ & h & main d'\oe uvre de l'exploitation, calculée sur son assolement 2017 observé \\
    $\sigma$ & sans & desserrement du plafond de travail, levier de scénario \\
    $r_c$ & unité/ha & taux par hectare d'un indicateur. Azote, \gls{ift}, \gls{ges}, eau,
      heures, euros de subvention \\
    $w_p$ & sans & poids de parcelle. Vaut 1, ou l'indicatrice d'irrigabilité pour l'eau \\
    $T$ & unité & niveau de la borne, fixé par le scénario \\
    \bottomrule
  \end{tabularx}
\end{table}

\section{Objectif et familles de contraintes}

\begin{table}[H]
  \centering
  \caption{Les cinq familles de contraintes du programme et les deux objectifs disponibles,
  avec l'équation du chapitre~\ref{ch:porter} qui donne la forme de chacune.}
  \label{tab:ann-familles}
  \scriptsize
  \setlength{\tabcolsep}{4pt}
  \begin{tabularx}{\textwidth}{@{}>{\raggedright\arraybackslash}p{30mm}
                               >{\raggedright\arraybackslash}p{20mm}X@{}}
    \toprule
    \textbf{Famille} & \textbf{Forme} & \textbf{Portée et remarque} \\
    \midrule
    Objectif Markowitz & \eqref{eq:objectif} & marge brute agrégée corrigée du risque. Objectif
      par défaut \\
    Objectif marge brute pure & \eqref{eq:objectif} sans le facteur de risque &
      activé par la seule politique de dérégulation totale. Sa colonne d'objectif est
      incomparable aux autres, et cette politique ne se juge que sur les grandeurs physiques \\
    \midrule
    Affectation & \eqref{eq:unicite} & au plus une culture par parcelle \\
    Règles par exploitation & \eqref{eq:mo} & plafond de travail, parts de surface, rotations \\
    Bornes de production & tonnes & le plancher de prairie et l'objectif de surface pâturée
      s'expriment en surface et non en tonnes \\
    Bornes d'indicateur & \eqref{eq:indicateur} & au territoire ou par zone. La forme zonale
      est celle qui écrit une directive nitrates vérifiée ferme par ferme \\
    Inertie & part de surface & exige qu'une part de la surface allouée reste dans son groupe
      de 2017. C'est ce qui distingue une transition d'une rupture \\
    \bottomrule
  \end{tabularx}
\end{table}

\begin{table}[H]
  \centering
  \caption{Composition du masque d'éligibilité. Chaque règle interdit son propre sous-ensemble
  de couples et le masque conserve l'union des interdictions.}
  \label{tab:ann-eligibilite}
  \scriptsize
  \setlength{\tabcolsep}{4pt}
  \begin{tabularx}{\textwidth}{@{}>{\raggedright\arraybackslash}p{34mm}X@{}}
    \toprule
    \textbf{Composante} & \textbf{Ce qu'elle interdit} \\
    \midrule
    Bornes numériques & altitude, pente, pluviométrie et surface de la parcelle hors de la
      plage admise par la culture \\
    Interdiction pure & une culture, sans condition \\
    Exigence d'irrigation & une culture sur une parcelle non irrigable \\
    Interdiction sur attribut & une culture sur une valeur d'attribut de la parcelle, par
      exemple un type de sol \\
    Seuil ou valeur exacte de risque & une culture au-delà d'un indicateur de risque, la règle
      chlordécone \\
    Interdiction régionale & une culture sur une île ou une région \\
    Verrou de friche & propre à l'étude de cas, fige les parcelles selon leur historique \\
    \midrule
    Conjonction et disjonction & une entrée unique conjugue ses conditions par un ET. Une
      disjonction s'écrit donc en plusieurs entrées, comme l'interdiction du melon, qui porte à
      la fois sur des types de sol et sur une île \\
    \bottomrule
  \end{tabularx}
\end{table}

\section{Indicateurs}

\begin{table}[H]
  \centering
  \caption{Indicateurs ajoutés au portage, en plus de l'azote, des \gls{ges} et de l'\gls{ift}
  hérités du \gls{gams}, et réserve de lecture attachée à chacun. Tous sont des taux par
  hectare appliqués à l'allocation après la résolution.}
  \label{tab:ann-indicateurs}
  \scriptsize
  \setlength{\tabcolsep}{4pt}
  \begin{tabularx}{\textwidth}{@{}p{30mm}X@{}}
    \toprule
    \textbf{Indicateur} & \textbf{Réserve de lecture} \\
    \midrule
    Phosphore, potasse & lus sur les noms des engrais NPK des itinéraires techniques. Lecture
      validée contre la colonne d'azote existante, qui la confirme au millième \\
    Eau & besoin brut des cultures, jamais un besoin net d'irrigation. La clé de répartition
      mensuelle des précipitations existe dans la documentation des données et n'était pas
      utilisée. C'est une correction immédiate, signalée en conclusion \\
    Carbone organique & seul indicateur à dépendre de la parcelle, par son type de sol, et non
      seulement de la culture. Il ne passe donc pas par l'aide au calcul commune aux autres \\
    Rpest & arbre flou de \textcite{tixier2007}, seul indicateur à dépendre du couple $(p,c)$.
      Il croise les produits phytosanitaires de la culture avec le ruissellement et le
      drainage de la parcelle \\
    Agroécologie, biologique & rapportés séparément. Les \gls{mae} donnent une définition
      sourcée de l'agroécologie, les opérations de fertilisation organique une définition du
      biologique. La canne récoltée en vert n'est pas biologique, et la prairie domine le
      biologique par itinéraire \\
    \midrule
    Score composite & pondère par famille. Sans cela les onze ratios d'autonomie alimentaire,
      quasi colinéaires, pèseraient onze fois les \gls{ges}, et le score mesurerait la finesse
      du découpage plutôt que la performance des politiques \\
    Indicateur borné par le scénario & exclu du score composite, sur détection automatique. Le
      noter reviendrait à féliciter une politique d'avoir respecté sa propre consigne \\
    \bottomrule
  \end{tabularx}
\end{table}

\begin{table}[H]
  \centering
  \caption{Deux modes de défaillance silencieux de la formulation, et la parade retenue pour
  chacun.}
  \label{tab:ann-pieges}
  \scriptsize
  \setlength{\tabcolsep}{4pt}
  \begin{tabularx}{\textwidth}{@{}>{\raggedright\arraybackslash}p{30mm}X@{}}
    \toprule
    \textbf{Piège} & \textbf{Comportement et parade} \\
    \midrule
    Le booléen trivial & lorsque la somme d'une contrainte ne rencontre aucun couple $(p,c)$
      éligible, elle vaut le \texttt{0} entier de Python et non une expression symbolique. La
      comparer à un seuil produit un booléen nu que le modeleur rejette. Toute contrainte dont
      l'ensemble de termes peut être vide rend donc explicitement \enquote{faisable} ou
      \enquote{infaisable} à la place \\
    Le registre muet & un nom de constructeur mal orthographié n'échoue qu'à l'interrogation du
      registre, et une famille dont le module n'est jamais importé est silencieusement absente.
      C'est le prix de l'indirection décrite au §~\ref{sec:architecture}, et il explique
      l'attachement du dépôt à ses garde-fous \\
    \bottomrule
  \end{tabularx}
\end{table}
